\documentclass[aspectratio=169,11pt]{beamer}

% ============================================================
% SISMID 2026 -- Day 2, 3:30 session, Part 1
% COVID-19 exercise (S. Yang + C. Djorno)
% Accessing digital traces + detecting exponential growth.
% Built on Mauricio's exercise 3 (MIGHTE-lab/SISMID25),
% covid_traces_WA.csv. Also previews the Day 3 capstone.
% Compile: pdflatex covid-exercise.tex  (run twice)
% ============================================================

\IfFileExists{beamerthememetropolis.sty}{
  \usetheme{metropolis}
  \metroset{progressbar=frametitle, sectionpage=progressbar, numbering=fraction}
}{
  \usetheme{Boadilla}
  \usecolortheme{seahorse}
  \setbeamertemplate{navigation symbols}{}
}

\usepackage[utf8]{inputenc}
\usepackage[T1]{fontenc}
\usepackage{tikz}
\usetikzlibrary{arrows.meta, positioning, fit, backgrounds, calc}
\usepackage{xcolor}
\usepackage{listings}
\usepackage{amsmath}

\definecolor{AccentA}{HTML}{2A6F97}
\definecolor{AccentB}{HTML}{C1666B}
\definecolor{AccentC}{HTML}{3B8A5B}
\setbeamercolor{frametitle}{bg=AccentA}
\setbeamercolor{progress bar}{fg=AccentB}

\definecolor{skbg}{rgb}{0.96,0.96,0.94}
\lstdefinestyle{prompt}{
  basicstyle=\ttfamily\scriptsize,
  backgroundcolor=\color{skbg},
  showstringspaces=false, breaklines=true,
  frame=single, framesep=4pt, columns=fullflexible, numbers=none,
}
\lstset{style=prompt}

\tikzset{
  cbox/.style={draw=AccentA, very thick, rounded corners, align=center,
    minimum height=1.0cm, inner sep=6pt, font=\small, fill=AccentA!6},
  sbox/.style={draw=AccentA, thick, rounded corners, align=center,
    minimum height=0.85cm, inner sep=4pt, font=\scriptsize, fill=AccentA!6},
  rbox/.style={draw=AccentB, thick, rounded corners, align=center,
    minimum height=0.85cm, inner sep=4pt, font=\scriptsize, fill=AccentB!8},
  gbox/.style={draw=AccentC, thick, rounded corners, align=center,
    minimum height=0.85cm, inner sep=4pt, font=\scriptsize, fill=AccentC!8},
  flow/.style={-{Stealth[length=3mm]}, very thick, draw=AccentB},
  flowq/.style={-{Stealth[length=2.4mm]}, thick, draw=AccentA!70},
}

\title{The COVID-19 Exercise}
\subtitle{Digital traces, and detecting when an outbreak is taking off}
\author{Shihao Yang \quad\textbullet\quad Candice Djorno}
\institute{SISMID 2026 \textbullet\ Day 2, 3:30}
\date{July 21, 2026}

\begin{document}

\begin{frame}
  \titlepage
\end{frame}

\begin{frame}{A different question}
  ARGO this afternoon asked \textbf{``how high is it?''} Early COVID asked a sharper one:
  \begin{block}{Is it taking off, right now?}
    Not the level, but the \textbf{onset}: is the curve in \textbf{exponential growth}? That
    is the earliest, most actionable signal an outbreak gives.
  \end{block}
  \vspace{0.3em}
  \begin{itemize}
    \item Data: \textbf{COVID-19 in Washington State}, case counts plus five digital traces.
    \item Same agent workflow; today the model is a growth detector, not a nowcaster.
    \item Continues Mauricio's exercise 3 (\texttt{covid\_traces\_WA.csv}).
  \end{itemize}
\end{frame}

\begin{frame}{Roadmap}
  \tableofcontents
\end{frame}

% ===============================================================
\section{Accessing the digital traces}

\begin{frame}{Five digital traces, one outbreak}
  \texttt{covid\_traces\_WA.csv} pairs the ground truth with the novel streams of the COVID
  era:
  \begin{itemize}
    \item \textbf{Ground truth:} \texttt{new\_cases} (positive tests).
    \item \textbf{Traces:} \texttt{upToDate} (clinician search), \texttt{cdc\_ili},
          \texttt{Twitter\_RelatedTweets}, \texttt{google\_fever}, \texttt{Kinsa}
          (smart-thermometer fever), \texttt{Cuebiq\_Mobility}.
  \end{itemize}
  \vspace{0.3em}
  \begin{center}
    \itshape The whole course in one file: many independent digital traces, one disease to
    track. (Pre-smoothed, so case counts are fractional.)
  \end{center}
\end{frame}

\begin{frame}[fragile]{Load and look first}
  \begin{block}{Prompt}
  \begin{lstlisting}
Load covid_traces_WA.csv (daily, Washington State). In a
6-panel plot, show new_cases and the five digital traces
over time. Which traces look like they rise before cases
(leading) and which after (lagging)?
  \end{lstlisting}
  \end{block}
  \begin{itemize}
    \item Before any algorithm, form the intuition by eye: where do outbreaks \emph{start},
          and do any traces move first?
    \item Leading traces are the prize: they are what an early-warning system watches.
  \end{itemize}
\end{frame}

% ===============================================================
\section{Detecting exponential growth}

\begin{frame}{What ``exponential growth'' means, operationally}
  Exponential growth is a \textbf{constant multiplicative factor} each step: today is roughly
  $\alpha$ times a recent baseline, with $\alpha > 1$.
  \begin{itemize}
    \item Over a sliding 11-day window, regress the \textbf{next 10 days on the previous 10}
          (through the origin, no intercept). The slope is $\alpha$.
    \item $\alpha > 1$ means the trace is \textbf{growing}; $\alpha < 1$ means it is
          receding.
  \end{itemize}
  \vspace{0.3em}
  \begin{center}
    \itshape Same idea as a log-linear fit: $\alpha \approx e^{r}$, so doubling time is
    $\ln 2 / \ln \alpha$. The slope is just an easy, robust way to read it.
  \end{center}
\end{frame}

\begin{frame}[fragile]{Compute the growth factor}
  \begin{block}{Prompt}
  \begin{lstlisting}
For new_cases, slide an 11-day window: in each window
regress the last 10 days on the previous 10 with no
intercept, and store the slope as alpha for that day.
Then make a 3-panel plot: cases, alpha over time with a
line at 1, and a 0/1 flag for alpha > 1.
  \end{lstlisting}
  \end{block}
  \begin{center}
    \itshape Now ``is it growing?'' is a number you can watch every day, not a judgment call.
  \end{center}
\end{frame}

% ===============================================================
\section{From growth to an outbreak flag}

\begin{frame}[fragile]{Turning the signal into a decision}
  A single day of $\alpha>1$ is noise. A \textbf{sustained} run is an outbreak.
  \begin{block}{Prompt}
  \begin{lstlisting}
Detect outbreak starts: when alpha > 1 for 10 consecutive
days, mark that 10th day as an outbreak start (only the
first such day). When alpha < 1 for 10 consecutive days,
the outbreak is over and a later run can start a new one.
Mark the detected starts on the case-count plot.
  \end{lstlisting}
  \end{block}
  \begin{center}
    \itshape A tiny state machine on top of $\alpha$: growing, active, over, repeat.
  \end{center}
\end{frame}

\begin{frame}{It finds the real Washington waves}
  Run on \texttt{new\_cases}, the detector flags starts that line up with the actual COVID
  waves:
  \begin{center}
  \footnotesize
  \texttt{2020-03-10} \ \textbullet\ \ \texttt{2020-06-06} \ \textbullet\ \
  \texttt{2020-09-27} \ \textbullet\ \ \texttt{2021-03-28 (Alpha)} \ \textbullet\ \
  \texttt{2021-07-21 (Delta)} \ \textbullet\ \ \texttt{2021-12-29 (Omicron)}
  \end{center}
  \vspace{0.4em}
  \begin{alertblock}{Verify, as always}
    Do the marks match where \emph{you} said outbreaks start by eye? A detector that fires
    late, or on every wobble, is worse than useless. Tune the window and threshold and see.
  \end{alertblock}
\end{frame}

% ===============================================================
\section{Multiple traces: early warning}

\begin{frame}[fragile]{Do the traces flag growth before cases?}
  The point of the digital traces: run the \emph{same} detector on each of them.
  \begin{block}{Prompt}
  \begin{lstlisting}
Repeat the alpha computation and outbreak detection for
each digital trace. Redo the 6-panel plot with each
trace's outbreak starts marked. For each trace, report how
many days before (or after) the case-count outbreak its
own outbreak fires.
  \end{lstlisting}
  \end{block}
  \begin{center}
    \itshape If enough traces flag growth ahead of cases, you can call an imminent outbreak.
    That is exactly Mauricio and Kogan's multi-trace early-warning idea.
  \end{center}
\end{frame}

% ===============================================================
\section{Wrap-up}

\begin{frame}{What you built}
  \begin{itemize}
    \item An \textbf{onset detector}: a growth factor $\alpha$ and a small state machine that
          calls the start of an outbreak, validated against the real WA waves.
    \item A \textbf{multi-trace} version that asks which digital streams lead, the seed of an
          early-warning system.
    \item It complements ARGO: ARGO says \emph{how high}, this says \emph{it is taking off}.
  \end{itemize}
  \vfill
  \begin{center}
    \itshape Two questions, two tools, both built from simple regressions on novel data.
  \end{center}
\end{frame}

\begin{frame}{Looking ahead to Day 3: bring your own problem}
  Tomorrow you build a situational-awareness product of your own. Start choosing now:
  \begin{itemize}
    \item \textbf{A disease:} flu, dengue, RSV, COVID, or something you work on.
    \item \textbf{A place:} your county, your country, Georgia for the capstone framing.
    \item \textbf{A question:} how high (ARGO), taking off (today), or where next (mobility)?
  \end{itemize}
  \vspace{0.3em}
  \begin{center}
    \itshape Every skill this week is a Lego brick. Tomorrow you pick which bricks your
    problem needs. (Mauricio's COVID research talk follows this exercise.)
  \end{center}
\end{frame}

\begin{frame}{Credits and sources}
  \footnotesize
  \begin{itemize}
    \item Exercise continues Mauricio Santillana's COVID exercise 3
          (\texttt{MIGHTE-lab/SISMID25}); data \texttt{covid\_traces\_WA.csv} (Washington
          State, daily, 2020--2021), pre-smoothed.
    \item Traces: UpToDate clinician search, CDC ILI, Twitter, Google ``fever'', Kinsa
          smart-thermometer fever, Cuebiq mobility.
    \item Multi-trace early warning: Kogan, Santillana et al., \emph{An early warning
          approach to monitor COVID-19 activity with multiple digital traces}, Science
          Advances 2021.
  \end{itemize}
\end{frame}

\begin{frame}[standout]
  Questions?\\[0.4em]
  \small Repo: \texttt{github.com/Shihao-Yang/sismid2026}
\end{frame}

\end{document}
